% CS HELP ===================================================== % 
% Instrukce k přílohám
% Do příloh závěrečné práce jsou umisťovány texty, obrázky, grafy, tabulky, které by přímo v~textu byly zbytečně detailní, ale mají být po ruce k~dokreslení východisek i~výsledků řešení.
% 
% Dále pak rozsáhlé/rozměrné grafické a~další objekty (např. výsledky měření, simulací, návrhy desek plošných spojů…), případně kompletní rozsáhlé a~souvislé výpisy zdrojových kódů programů a~skriptů, které vzhledem ke svému rozsahu nebo významu (nejsou zcela podstatné pro výklad provedený v~dané kapitole/odstavci) není vhodné umisťovat do vlastních kapitol/textu práce, protože by způsobily značný nárůst počtu jejich stran a~došlo by ke snížení přehlednosti práce.
% 
% Pokud však autor práce v~určité kapitole vysvětluje/popisuje např. princip funkce navrženého elektronického obvodu, mělo by být i~rozsáhlejší schéma zapojení (nebo alespoň jeho výřez) umístěno přímo v~textu práce v~dané kapitole, kde je popis proveden z~důvodu lepší čitelnosti. Toto se týká i~dalších rozměrnějších (např. grafických) objektů, jejichž význam je zcela zásadní pro daný text v~dílčích kapitolách práce, tzn. není nutné mít přehnané obavy z~umístění rozměrnějších grafických objektů mimo přílohy, pokud toto např. zvýší přehlednost práce a~její srozumitelnost.
% 
% Všechny objekty umístěné v~přílohách musí obsahovat titulek s~příslušným indexováním nebo číslováním (případně také stručný popis), dílčí části příloh by také měly být vhodně logicky rozčleněny včetně uvedení vhodných nadpisů jednotlivých částí/kapitol příloh. Logika rozčlenění a~popisu by měla být ve shodě s~vlastním textem v~jednotlivých kapitolách práce, ze kterých je odkazováno na příslušnou část přílohy (je vhodné využívat křížové odkazy z~vlastního textu práce do příloh).
% 
% Všechny grafické výstupy/objekty vložené do příloh musí být v~maximální možné grafické kvalitě, tzn. je vhodné používat výhradně vektorové formáty (*.eps, *.ps, *.hpgl, *.svg, *.wmf, *.emf  atd.) pro export z~návrhových systémů pro desky plošných spojů, Matlabu atd., rastrové formáty (*.jpg, *.bmp, *.gif, *.tif  atd.) jsou pro tyto účely nevhodné.
% 
% Pokud je v~příloze např. ve schématu uveden text, hodnoty součástek a~další popisky, musí být bez problému čitelné. Pokud toto není možné zajistit, je nutné přílohu vytisknout na větší formát papíru, případně ji rozdělit na více stran. Ve výjimečných případech, schválených děkanem FEL, může zvláštní část příloh obsahovat utajované výsledky, které by neměly být zveřejněné spolu s~celou prací.
% ============================================================= %


\chapter{Vzorová příloha}

\lipsum[3-5]


